% META: {"id": "TSPE-PROBABILITES-QCM", "chapitre": "TSPE-PROBABILITES", "type_objet": "qcm", "genere_depuis": "chapitres/TSPE-PROBABILITES/qcm/TSPE-PROBABILITES-QCM.json", "status": "generated"}
% Fichier genere par scripts/build_qcm_tex.py — ne pas editer a la main.

\section*{\textcolor{chapcolor}{\MakeUppercase{QCM d'auto-évaluation -- Probabilités}}}

\begin{center}
\textit{Pour chaque question, une seule réponse est exacte.}
\end{center}

\begin{enumerate}

\item[] \textbf{Capacité C5}

\item \textbf{[Q1]} Si X suit B(n,p), alors P(X=k) vaut :
  \begin{enumerate}[label=\Alph*.]
    \item C(n,k) p\textasciicircum{}k (1-p)\textasciicircum{}(n-k)
    \item p\textasciicircum{}k (1-p)\textasciicircum{}(n-k)
    \item C(n,k) p\textasciicircum{}k
    \item np\textasciicircum{}k
  \end{enumerate}

\item[] \textbf{Capacité C2}

\item \textbf{[Q2]} Un schema de Bernoulli suppose que les epreuves sont :
  \begin{enumerate}[label=\Alph*.]
    \item dependantes et identiques
    \item indépendantes et identiques
    \item indépendantes et différentes
    \item toutes de probabilité 1/2
  \end{enumerate}

\item[] \textbf{Capacité C7}

\item \textbf{[Q3]} La linéarité de l'espérance E(X+Y)=E(X)+E(Y) est valable :
  \begin{enumerate}[label=\Alph*.]
    \item toujours, sans condition d'independence
    \item uniquement si X et Y sont indépendantes
    \item uniquement si X=Y
    \item jamais
  \end{enumerate}

\item[] \textbf{Capacité C8}

\item \textbf{[Q4]} L'additivite de la variance V(X+Y)=V(X)+V(Y) est garantie lorsque :
  \begin{enumerate}[label=\Alph*.]
    \item que X=Y
    \item que X et Y soient indépendantes
    \item rien de particulier
    \item que E(X)=E(Y)
  \end{enumerate}

\item[] \textbf{Capacité C9}

\item \textbf{[Q5]} Pour X \textasciitilde{} B(n,p), V(X) vaut :
  \begin{enumerate}[label=\Alph*.]
    \item np
    \item n(1-p)
    \item np(1-p)
    \item p(1-p)
  \end{enumerate}

\item[] \textbf{Capacité C10}

\item \textbf{[Q6]} L'inégalité de Bienayme-Tchebychev majoré :
  \begin{enumerate}[label=\Alph*.]
    \item P(Y=mu)
    \item E(Y)
    \item V(Y)
    \item P(|Y-mu|>=delta)
  \end{enumerate}

\end{enumerate}

% NEXUS-QCM-TEACHER-ONLY-BEGIN
\ifnxVersionProfesseur
\clearpage
\section*{Clé de correction — réservée au professeur}

\begin{center}
\begin{tabular}{lll}
\hline
Question & Capacité & Réponse exacte \\
\hline
Q1 & C5 & \textbf{A} \\
Q2 & C2 & \textbf{B} \\
Q3 & C7 & \textbf{A} \\
Q4 & C8 & \textbf{B} \\
Q5 & C9 & \textbf{C} \\
Q6 & C10 & \textbf{D} \\
\hline
\end{tabular}
\end{center}

\subsection*{Diagnostic des réponses erronées}

\begin{itemize}
  \item \textbf{Q1}
  \begin{itemize}
    \item \textbf{B} — Cette expression donne la probabilité d'une configuration précise de $k$ succès ; elle oublie les $C(n,k)$ positions possibles des succès. \emph{Renvoi : C5.}
    \item \textbf{C} — Le facteur $(1-p)^{n-k}$ manque : les $n-k$ echecs de chaque configuration doivent eux aussi contribuer a sa probabilité. \emph{Renvoi : C5.}
    \item \textbf{D} — Le produit $np^k$ ne compte ni les choix des $k$ succès ni la probabilité des $n-k$ echecs ; ce n'est pas une probabilité binomiale. \emph{Renvoi : C5.}
  \end{itemize}
  \item \textbf{Q2}
  \begin{itemize}
    \item \textbf{A} — La dépendance contredit la définition du schema : la probabilité d'un succès doit rester $p$ d'une epreuve a l'autre.
    \item \textbf{C} — Des epreuves différentes peuvent avoir des probabilités de succès distinctes ; un schema de Bernoulli repete la même epreuve.
    \item \textbf{D} — La probabilité de succès est un réel $p$ quelconque de $[0,1]$ ; la valeur $1/2$ n'est imposee que pour une piece equilibree.
  \end{itemize}
  \item \textbf{Q3}
  \begin{itemize}
    \item \textbf{B} — L'indépendance est requise pour additionner les variances, pas pour la linéarité de l'espérance de variables integrables.
    \item \textbf{C} — L'égalité $X=Y$ n'est pas nécessaire : la somme des esperances vaut l'espérance de la somme pour deux variables quelconques.
    \item \textbf{D} — La propriété est toujours valable ; par exemple pour deux constantes $X=1$ et $Y=2$, les deux membres valent $3$.
  \end{itemize}
  \item \textbf{Q4}
  \begin{itemize}
    \item \textbf{A} — Si $X=Y$, alors $V(X+Y)=V(2X)=4V(X)$, tandis que $V(X)+V(Y)=2V(X)$ : l'égalité n'est donc pas garantie.
    \item \textbf{C} — Sans hypothèse, la formule comporte le terme $2\operatorname{Cov}(X,Y)$ : l'indépendance permet precisement d'annuler cette covariance.
    \item \textbf{D} — L'égalité des esperances ne renseigne pas sur la covariance de $X$ et $Y$ : elle ne suffit donc pas a rendre les variances additives.
  \end{itemize}
  \item \textbf{Q5}
  \begin{itemize}
    \item \textbf{A} — Le produit $np$ est l'espérance de la loi binomiale ; la variance comporte aussi le facteur de dispersion $1-p$.
    \item \textbf{B} — L'expression $n(1-p)$ oublie le facteur $p$ associe aux succès ; chaque indicatrice de Bernoulli a pour variance $p(1-p)$.
    \item \textbf{D} — Le produit $p(1-p)$ est la variance d'une seule epreuve de Bernoulli ; la somme de $n$ epreuves indépendantes apporte le facteur $n$.
  \end{itemize}
  \item \textbf{Q6}
  \begin{itemize}
    \item \textbf{A} — Bienayme-Tchebychev concerne l'événement $|Y-\mu|\geq\delta$, pas la probabilité ponctuelle que $Y$ soit exactement égal a $\mu$.
    \item \textbf{B} — L'espérance $E(Y)$ est le centre de l'événement de deviation ; l'inégalité ne cherche pas a majorer cette valeur.
    \item \textbf{C} — La variance intervient dans le majorant $V(Y)/\delta^2$, mais la quantité majorée est la probabilité d'un ecart a la moyenne.
  \end{itemize}
\end{itemize}
\fi
% NEXUS-QCM-TEACHER-ONLY-END
